%% text-sharing.tex -- shared closing section
\subsubsection{Building and sharing}

The package consists of \texttt{florilegium.sty}, the engine-specific
\texttt{.def} hosts and the Lua sources beside them. Drawing and composition live
in \texttt{ink}, \texttt{flora}, \texttt{authored} and \texttt{grow}; the other
modules supply layout settings, obstacle routing, decorated initials, botanical
lettering and the host interface. A replaceable drawing sink lets the same paths
be written inside Lua\LaTeX\ or to a standalone vector PDF.

Lua\LaTeX\ can render a simple page in one pass. Xe\LaTeX\ records the drawings
and text marks, an external \texttt{texlua} pass renders them, and the next TeX
pass places the vector pages. The build command manages these passes, reference
changes and aging without shell escape. A source kit includes a book template,
a four-page design guide, the license and a file-hash manifest; its extracted
template is checked with both engines.

\begin{remark}
Two mistakes cost time and are worth recording. A literal written into the page bypasses
TikZ, so it also bypasses TikZ's transformations: a primitive in a shifted or rotated
\texttt{scope} must read pgf's current matrix itself (it is in
\verb|\pgf@pt@aa| to \verb|\pgf@pt@bb| and the dimensions \verb|\pgf@pt@x|,
\verb|\pgf@pt@y|) and apply it in Lua. And a Lua function that returns two values, placed
in the middle of a table constructor such as \verb|{x0, y0, f(t), x1, y1}|, contributes only
its first value; the rest of the table shifts, and the error surfaces three functions
later as arithmetic on \texttt{nil}.
\end{remark}
